\begin{abstract}
Tactile sensing is emerging as an indispensable modality for dexterous robotic manipulation, driven by the convergence of low-cost vision-based sensors, open-source robot hands, and large-scale foundation models. This survey provides a comprehensive review of tactile sensing for robot hands, covering six tightly coupled dimensions: (1) tactile sensor technologies, from piezoresistive transduction to the GelSight--DIGIT--Digit~360 lineage of vision-based optical sensors; (2) robot hand design, including the open-source revolution (LEAP Hand at \$2K, ORCA, ISyHand) and intelligent mechanisms such as underactuation, variable stiffness actuators, and active surfaces; (3) tactile data representation and collection, anchored by the six-structure taxonomy of Albini et al.\ and self-supervised foundation models (Sparsh, UniTouch); (4) learning-based manipulation, spanning Diffusion Policy, reinforcement learning, and force-informed approaches (DexForce, ForceVLA); (5) vision-language-action (VLA) models from RT-2 to pi0 and Gemini Robotics, with emerging tactile integration; and (6) human-to-robot transfer via embodiment retargeting across kinematic, visual, and tactile domains. We synthesize 146 papers (2020--2026), identify the top~10 open challenges---led by the tactile sim-to-real gap and cross-embodiment tactile transfer---and propose a mechanism--tactile--learning triangle as a unifying research direction.
\end{abstract}
